\chapter{NỘI DUNG THỰC TẬP}

\section{Phạm vi và phương pháp tìm hiểu}

Trong thời gian thực tập, dưới sự hướng dẫn của anh \CompanySupervisor{}, em tập
trung tìm hiểu hoạt động cung ứng linh kiện điện tử và thiết bị tự động hóa tại
\CompanyName{}. Em chọn đi vào ba vấn đề chính: thứ nhất là mối liên hệ giữa các
bộ phận tham gia chuỗi cung ứng; thứ hai là cách chuyển yêu cầu của khách hàng
thành tiêu chí chọn thiết bị; và thứ ba là cách kiểm tra một mã hàng cụ thể bằng
tài liệu của nhà sản xuất.

Nguồn thông tin chủ yếu là website giới thiệu, trang bán hàng công khai của
TECHPAL và tài liệu kỹ thuật (datasheet) từ các hãng sản xuất. Em chỉ sử dụng
những gì có thể tra cứu và kiểm chứng được, còn quy trình nội bộ, dữ liệu kho
hay kết quả đo kiểm thì em không có quyền truy cập. Các phân tích ứng dụng trong
chương này mới ở mức phương án trên giấy, chưa phải sản phẩm đã làm ra.

Em thực hiện theo bốn bước:

\begin{enumerate}[label=\arabic*)]
  \item Tìm các nhóm sản phẩm và bộ phận liên quan từ nguồn công khai;
  \item Chọn vài mã đại diện có trong danh mục của TECHPAL;
  \item Đối chiếu mô tả trên trang bán hàng với tài liệu chính thức của hãng;
  \item Tổng hợp lại thành quy trình lựa chọn, kèm một bài toán minh họa.
\end{enumerate}

\section{Các vị trí liên quan trong chuỗi cung ứng}

Trên trang giới thiệu của TECHPAL có đề cập đến ba bộ phận trực tiếp tham gia
cung ứng: Sale, Mua hàng Quốc tế và Logistics. Ba bộ phận
này phụ trách các giai đoạn khác nhau nhưng đều làm việc trên cùng một đơn hàng,
cùng mã linh kiện, số lượng, yêu cầu chất lượng và deadline giao hàng
\cite{techpal-about}.

\begin{table}[H]
\centering
\caption{Vai trò của các bộ phận trong chuỗi cung ứng linh kiện}
\label{tab:chapter2-roles}
\small
\renewcommand{\arraystretch}{1.1}
\begin{tabularx}{\textwidth}{>{\bfseries\raggedright\arraybackslash}p{3.1cm} >{\raggedright\arraybackslash}X >{\raggedright\arraybackslash}X}
\toprule
Bộ phận & Nội dung công việc liên quan & Kiến thức, kỹ năng cần thiết \\
\midrule
Sale & Tiếp nhận nhu cầu, làm rõ mã hàng, số lượng, ứng dụng và thời hạn; trao đổi phương án với khách hàng & Giao tiếp kỹ thuật, đọc mã linh kiện, hiểu thông số cơ bản và lập báo giá \\
Mua hàng Quốc tế & Tìm nguồn, kiểm tra khả năng cung ứng, quy cách đóng gói và chứng từ; làm việc với nhà cung cấp & Tiếng Anh chuyên ngành, tra cứu chéo, đánh giá nguồn và quản lý thay đổi sản phẩm \\
Logistics & Theo dõi vận chuyển, tiếp nhận, bảo quản và bàn giao; kiểm tra sự phù hợp giữa hàng hóa với chứng từ & Quản lý mã hàng, đóng gói, điều kiện lưu kho, ESD và truy xuất lô hàng \\
\bottomrule
\end{tabularx}
\end{table}

Bảng~\ref{tab:chapter2-roles} là do em tự tổng hợp theo tính chất công việc, không
phải bản mô tả nhân sự chính thức của công ty. Tuy nhiên, một điểm em thấy rõ là
với linh kiện điện tử, chỉ cần sai một hậu tố trong mã đặt hàng là có thể nhận
nhầm kiểu vỏ, dải nhiệt độ hay quy cách đóng gói. Cho nên thông tin kỹ thuật
cần phải được giữ nhất quán khi chuyển từ bộ phận này sang bộ phận khác.

Từ thông tin về dịch vụ đặt hàng qua châu Âu -- Mỹ, DigiKey, Mouser và Trung
Quốc trên trang TECHPAL \cite{techpal-about}, em vẽ lại luồng xử lý một
yêu cầu cung ứng như Hình~\ref{fig:supply-flow}. Sơ đồ này chỉ mang tính khái
quát để phục vụ học tập, không phải quy trình nội bộ mà công ty đang áp dụng.

\begin{figure}[H]
\centering
\begin{tikzpicture}[
  block/.style={draw, rounded corners=2pt, align=center, minimum height=1.25cm,
                text width=2.35cm, font=\small},
  line/.style={-{Latex[length=2.3mm]}, thick},
  node distance=0.45cm
]
  \node[block] (need) {Tiếp nhận\\yêu cầu};
  \node[block, right=of need] (normalize) {Chuẩn hóa\\mã và thông số};
  \node[block, right=of normalize] (source) {Tìm nguồn và\\đối chiếu};
  \node[block, right=of source] (order) {Đặt hàng,\\vận chuyển};
  \node[block, right=of order] (deliver) {Tiếp nhận và\\bàn giao};
  \draw[line] (need) -- (normalize);
  \draw[line] (normalize) -- (source);
  \draw[line] (source) -- (order);
  \draw[line] (order) -- (deliver);
\end{tikzpicture}
\caption{Luồng xử lý yêu cầu cung ứng (mô hình hóa từ thông tin công khai)}
\label{fig:supply-flow}
\end{figure}

\section{Khảo sát danh mục sản phẩm của TECHPAL}

Danh mục trực tuyến của TECHPAL chia sản phẩm thành bốn nhóm lớn: linh kiện điện
tử; thiết bị tự động hóa công nghiệp; vật tư phụ trợ công nghiệp; và thiết bị,
vật tư phòng sạch. Cách phân nhóm này
giúp thu hẹp phạm vi tìm kiếm, nhưng chỉ nhìn vào tên nhóm thì chưa biết được
sản phẩm có hợp với ứng dụng cụ thể hay không \cite{techpal-components}.

\begin{table}[H]
\centering
\caption{Các nhóm cung ứng chính trên danh mục TECHPAL}
\label{tab:techpal-categories}
\small
\renewcommand{\arraystretch}{1.1}
\begin{tabularx}{\textwidth}{>{\bfseries\raggedright\arraybackslash}p{4.1cm} >{\raggedright\arraybackslash}X >{\raggedright\arraybackslash}X}
\toprule
Nhóm cung ứng & Thành phần tiêu biểu quan sát được & Thông tin cần kiểm tra \\
\midrule
Linh kiện điện tử & IC, vi điều khiển, linh kiện công suất, đầu nối, rơ-le, màn hình & Mã đầy đủ, kiểu vỏ, điện áp, dòng, nhiệt độ và trạng thái vòng đời \\
Thiết bị tự động hóa & Mô-đun PLC/I/O, cảm biến, encoder, bộ nguồn, thiết bị ghi dữ liệu & Kiến trúc hệ thống, số kênh, loại tín hiệu, giao thức và khả năng mở rộng \\
Vật tư phụ trợ công nghiệp & Đầu nối, dây dẫn, phụ kiện lắp đặt và vật tư bảo trì & Kích thước, vật liệu, điều kiện môi trường và tương thích cơ khí \\
Thiết bị, vật tư phòng sạch & Dụng cụ và vật tư phục vụ môi trường kiểm soát bụi bẩn & Cấp sạch, vật liệu, quy cách đóng gói và yêu cầu sử dụng \\
\bottomrule
\end{tabularx}
\end{table}

Để không dừng lại ở mô tả chung chung, em chọn bốn mã hàng có thật trong danh
mục TECHPAL \cite{techpal-components} để tra cứu kỹ hơn. Bốn mã này đại diện cho bốn loại thiết bị khác
nhau: IC truyền thông, mô-đun I/O công nghiệp, thiết bị kết nối I/O phân tán và
bộ ghi dữ liệu.

\begin{table}[H]
\centering
\caption{Một số mã hàng đại diện trong danh mục TECHPAL}
\label{tab:representative-products}
\small
\renewcommand{\arraystretch}{1.1}
\begin{tabularx}{\textwidth}{>{\ttfamily\raggedright\arraybackslash}p{4.1cm} >{\raggedright\arraybackslash}p{2.8cm} >{\raggedright\arraybackslash}X}
\toprule
\normalfont\textbf{Mã sản phẩm} & \textbf{Nhóm chức năng} & \textbf{Mô tả dùng cho bước tra cứu} \\
\midrule
MAX487CSA+T & IC truyền thông & Bộ thu phát RS-485/RS-422 công suất thấp, vỏ SOIC 8 chân \\
BMXAMI0410 & Mô-đun vào tương tự & Mô-đun Modicon X80 có bốn ngõ vào tương tự tốc độ cao \\
6ES7153-1AA03-0XB0 & Giao tiếp I/O phân tán & Mô-đun kết nối SIMATIC DP IM 153-1 cho ET 200M \\
CR1000X & Đo lường và lưu trữ & Bộ ghi dữ liệu có khả năng đo, xử lý, lưu trữ và truyền thông \\
\bottomrule
\end{tabularx}
\end{table}

\section{Quy trình đọc và lựa chọn thiết bị}

\subsection{Chuẩn hóa yêu cầu đầu vào}

Nếu khách hàng chỉ nói ``cần IC RS-485'' hay ``cần mô-đun analog'' thì chưa đủ
thông tin để đặt hàng. Yêu cầu như vậy cần được cụ thể hóa thành: chức năng gì,
điện áp bao nhiêu, mấy kênh, loại tín hiệu, tốc độ, giao thức, nhiệt độ làm
việc, kiểu lắp, số lượng và khi nào cần. Nếu đây là mã thay thế cho linh kiện
cũ thì phải có thêm mã gốc và thông tin về bo mạch đang sử dụng.

\subsection{Kiểm tra mã và tài liệu kỹ thuật}

Mã hàng cần đọc đầy đủ kể cả hậu tố -- điểm này ban đầu em cũng hay bỏ qua.
Để xác nhận thông số, nguồn đáng tin nhất là trang sản phẩm và datasheet của nhà
sản xuất. Mô tả trên website nhà cung ứng giúp định vị nhanh sản phẩm, nhưng
không nên dựa vào đó để kết luận tương thích hay không. Các bước kiểm tra em đề
xuất gồm:

\begin{enumerate}[label=\arabic*)]
  \item Xác nhận nhà sản xuất và mã đặt hàng đầy đủ;
  \item Đối chiếu chức năng, sơ đồ chân hoặc kiến trúc mô-đun;
  \item Kiểm tra giới hạn điện, nhiệt và điều kiện làm việc khuyến nghị;
  \item Xem kiểu vỏ, kích thước, đầu nối và cách lắp đặt;
  \item Xác định trạng thái vòng đời: còn sản xuất hay chỉ còn dạng phụ tùng;
  \item Lưu lại link tài liệu, ngày tra cứu và kết luận.
\end{enumerate}

Kết quả tra cứu nên ghi lại dưới dạng bảng kiểm để sau này người khác có thể
xem lại được nguồn dữ liệu và lý do chọn. Bảng~\ref{tab:technical-checklist}
liệt kê những câu hỏi tối thiểu cần trả lời trước khi đề xuất mua một mã hàng.

\begin{table}[H]
\centering
\caption{Bảng kiểm kỹ thuật trước khi lựa chọn mã hàng}
\label{tab:technical-checklist}
\footnotesize
\renewcommand{\arraystretch}{1.0}
\begin{tabularx}{\textwidth}{>{\bfseries\raggedright\arraybackslash}p{3.4cm} >{\raggedright\arraybackslash}X >{\raggedright\arraybackslash}X}
\toprule
Nội dung kiểm tra & Câu hỏi cần trả lời & Rủi ro nếu bỏ sót \\
\midrule
Định danh & Mã, hậu tố và hãng đã chính xác chưa? & Nhận sai biến thể hoặc sai quy cách đóng gói \\
Thông số điện & Điện áp, dòng, mức logic và công suất có phù hợp? & Hỏng linh kiện hoặc hệ thống hoạt động không ổn định \\
Giao tiếp & Chuẩn vật lý, giao thức và tốc độ có tương thích? & Không trao đổi được dữ liệu \\
Cơ khí & Vỏ, chân, đầu nối và kích thước lắp có khớp? & Không lắp được hoặc phải sửa thiết kế \\
Môi trường & Nhiệt độ, độ ẩm, rung và cấp bảo vệ có đáp ứng? & Suy giảm độ tin cậy khi vận hành \\
Vòng đời & Sản phẩm còn sản xuất hay chỉ còn dạng phụ tùng? & Khó mua bổ sung và tăng chi phí bảo trì \\
Chứng từ & Datasheet, nguồn gốc, lô hàng và chứng nhận có đủ? & Khó truy xuất và nghiệm thu \\
\bottomrule
\end{tabularx}
\end{table}

\subsection{Đánh giá khả năng thay thế}

Hai sản phẩm có cùng tên chức năng chưa chắc đã thay thế được cho nhau. Với IC
thì phải so chân, mức điện, tốc độ, tải bus, vỏ và dải nhiệt độ. Với mô-đun
công nghiệp thì còn phải xem rack, bộ điều khiển đi kèm, phiên bản firmware, cấu
hình I/O và chứng nhận. Đặc biệt khi sản phẩm đã ngừng sản xuất (chỉ còn dạng
phụ tùng), chi phí thay thế phải tính cả việc chuyển đổi hệ thống chứ không chỉ
nhìn giá mô-đun.

\section{Phân tích một số mã sản phẩm đại diện}

\subsection{Bộ thu phát MAX487CSA+T}

Trên trang TECHPAL, MAX487CSA+T được liệt kê trong nhóm linh kiện điện tử, mô
tả là bộ thu phát bán song công, vỏ SOIC 8 chân \cite{techpal-max487}. Em tra
trên trang Analog Devices thì thấy MAX487 là IC thu phát RS-485/RS-422 nguồn đơn
5~V, phiên bản giới hạn tốc độ sườn, hỗ trợ tới 250~kbps. Đầu vào thu có tải
$1/4$ đơn vị nên về lý thuyết có thể nối tối đa 128 nút trên một bus
\cite{adi-max487}.

Nếu dùng MAX487CSA+T cho giao tiếp nối tiếp trong thiết kế nhúng thì cần kiểm
tra thêm mấy thứ: điện áp logic phía vi điều khiển có tương thích 5~V không,
đường truyền có cần điện trở kết thúc không, bus cần phân cực ra sao, có cần bảo
vệ ESD không, và nếu hai đầu chênh lệch điện thế đất thì có cần cách ly không.
Ngoài ra, hậu tố ``T'' liên quan tới quy cách đóng gói (tape and reel) nên khi
đặt mua không được bỏ.

\subsection{Mô-đun vào tương tự BMXAMI0410}

Mã BMXAMI0410 trên TECHPAL thuộc nhóm thiết bị tự động hóa. Theo Schneider
Electric, đây là mô-đun Modicon X80 với bốn ngõ vào analog tốc độ cao
\cite{schneider-bmxami0410}. Khác với một IC rời, mô-đun này không thể dùng
đơn lẻ mà phụ thuộc vào cả hệ thống: rack nào tương thích, CPU PLC nào, cảm
biến ra kiểu tín hiệu gì, đấu dây thế nào, phần mềm cấu hình ra sao.

Điều em rút ra ở đây là trong cung ứng tự động hóa, số kênh hay loại tín hiệu
mới chỉ là điều kiện cần. Muốn chốt mã thì phải xem thêm tài liệu tương thích
hệ thống và phụ kiện đầu nối. Giả sử khách hàng chỉ nói ``cho mô-đun analog bốn
kênh'', bên kỹ thuật hoặc Sale cần hỏi thêm đang dùng PLC/rack gì và đo tín
hiệu dải nào.

\subsection{Mô-đun Siemens 6ES7153-1AA03-0XB0}

TECHPAL có bán mã 6ES7153-1AA03-0XB0. Tra trên trang Siemens, em thấy đây là
mô-đun kết nối SIMATIC DP IM~153-1 dùng cho trạm ET~200M, hỗ trợ tối đa tám
mô-đun S7-300. Điểm đáng chú ý là mã này hiện tại Siemens phân loại là
\textit{phụ tùng} (spare part only) \cite{siemens-im153}.

Trạng thái phụ tùng ảnh hưởng rất lớn đến quyết định mua. Nếu nhà máy đang chạy
hệ thống cũ và cần thay thế để giảm thời gian dừng máy thì vẫn phải mua mã này.
Nhưng nếu làm dự án mới thì nên cân nhắc nền tảng khác có vòng đời dài hơn. Em
nghĩ đây là ví dụ rõ nhất cho thấy kiểm tra vòng đời quan trọng không kém kiểm
tra thông số kỹ thuật.

\subsection{Bộ ghi dữ liệu CR1000X}

TECHPAL liệt kê CR1000X của Campbell Scientific trong nhóm thiết bị tự động hóa,
mô tả là bộ lưu trữ xử lý dữ liệu. Tài liệu hãng cho
biết CR1000X có 16 đầu vào analog đơn (hoặc 8 vi sai), ADC 24 bit, và tích hợp
sẵn khả năng đo cảm biến, xử lý, lưu trữ lẫn truyền dữ liệu
\cite{techpal-cr1000x,campbell-cr1000x}.

CR1000X khá phù hợp cho các bài toán thu thập dữ liệu môi trường hay trạm đo
từ xa. Tuy vậy, không phải cứ thấy 16 kênh là đủ. Người chọn thiết bị phải xem
kỹ: loại cảm biến, dải đo, độ chính xác cần thiết, chu kỳ lấy mẫu, bộ nhớ, cổng
giao tiếp, nguồn cấp và điều kiện ngoài trời hay trong nhà. Hãng cũng khuyến cáo
nên xác nhận cấu hình trước khi đặt mua.

\section{Bài toán khảo sát: thu thập và truyền dữ liệu nhiệt độ}

\subsection{Yêu cầu bài toán}

Để áp dụng thử quy trình đã nêu, em tự đặt ra một bài toán: thiết kế phương án
thu thập nhiệt độ trong tủ kỹ thuật. Hệ thống cần đo được nhiệt độ, lưu giá trị
theo thời gian, cảnh báo khi vượt ngưỡng và truyền dữ liệu về máy tính giám sát.
Cần nhấn mạnh rằng đây chỉ là phương án trên giấy, dùng để so sánh các nhóm
thiết bị trong danh mục TECHPAL, không phải hệ thống đã lắp đặt hay thử nghiệm.

\begin{table}[H]
\centering
\caption{Yêu cầu đầu vào của bài toán khảo sát}
\label{tab:temperature-requirements}
\begin{tabularx}{\textwidth}{>{\bfseries\raggedright\arraybackslash}p{4.3cm} >{\raggedright\arraybackslash}X}
\toprule
Đại lượng đo & Nhiệt độ không khí trong tủ kỹ thuật \\
Số điểm đo & Từ một điểm, có khả năng mở rộng \\
Chu kỳ lưu & Có thể cấu hình theo tốc độ biến đổi nhiệt \\
Chức năng & Hiển thị, lưu dữ liệu và cảnh báo vượt ngưỡng \\
Truyền thông & Kết nối tới máy tính hoặc mạng giám sát \\
Nguồn & Chọn theo thiết bị, có bảo vệ và dự phòng phù hợp \\
Giới hạn báo cáo & Phân tích kiến trúc, chưa chế tạo và chưa có số liệu đo \\
\bottomrule
\end{tabularx}
\end{table}

\subsection{Kiến trúc đề xuất}

Tùy vào hệ thống hiện có của khách hàng, phần thu thập dữ liệu có thể là bộ ghi
chuyên dụng (như CR1000X), mô-đun analog trong hệ PLC (như BMXAMI0410), hoặc một
bo mạch nhúng tự thiết kế dùng IC RS-485 (như MAX487). Ba hướng này không thay
thế nhau trực tiếp mà phù hợp với các mức độ tích hợp khác nhau.

\begin{figure}[H]
\centering
\begin{tikzpicture}[
  block/.style={draw, rounded corners=2pt, align=center, minimum height=1.1cm,
                text width=2.30cm, font=\small},
  line/.style={-{Latex[length=2.3mm]}, thick},
  node distance=0.35cm
]
  \node[block] (sensor) {Cảm biến\\nhiệt độ};
  \node[block, right=of sensor] (condition) {Đấu nối và\\điều hòa tín hiệu};
  \node[block, right=of condition] (acquire) {Khối thu thập\\và xử lý};
  \node[block, right=of acquire] (storage) {Lưu trữ và\\cảnh báo};
  \node[block, right=of storage] (monitor) {Máy tính hoặc\\mạng giám sát};
  \draw[line] (sensor) -- (condition);
  \draw[line] (condition) -- (acquire);
  \draw[line] (acquire) -- (storage);
  \draw[line] (storage) -- (monitor);
\end{tikzpicture}
\caption{Kiến trúc khái niệm của hệ thống thu thập nhiệt độ}
\label{fig:temperature-architecture}
\end{figure}

\begin{table}[H]
\centering
\caption{Đánh giá vai trò của các mã sản phẩm trong phương án}
\label{tab:product-role-evaluation}
\small
\renewcommand{\arraystretch}{1.1}
\begin{tabularx}{\textwidth}{>{\ttfamily\raggedright\arraybackslash}p{4.1cm} >{\raggedright\arraybackslash}X >{\raggedright\arraybackslash}X}
\toprule
\normalfont\textbf{Mã sản phẩm} & \textbf{Vai trò có thể xem xét} & \textbf{Điều kiện cần xác minh} \\
\midrule
CR1000X & Thu thập, xử lý và lưu dữ liệu trong một thiết bị chuyên dụng & Loại cảm biến, dải đo, chu kỳ lấy mẫu, nguồn và giao tiếp \\
BMXAMI0410 & Nhận tín hiệu analog trong hệ PLC Modicon X80 & Rack, CPU, đầu nối, loại tín hiệu và cấu hình phần mềm \\
MAX487CSA+T & Tạo giao tiếp RS-485 cho phương án bộ xử lý nhúng & Nguồn 5~V, mức logic, topology bus, kết thúc và bảo vệ \\
6ES7153-1AA03-0XB0 & Duy trì kết nối I/O phân tán trong hệ ET 200M hiện hữu & Cấu hình trạm, tương thích S7-300 và trạng thái phụ tùng \\
\bottomrule
\end{tabularx}
\end{table}

Thực tế thì chưa thể chốt phương án nào cho đến khi biết rõ khách hàng dùng cảm
biến gì, đang có sẵn nền tảng điều khiển nào và yêu cầu vận hành cụ thể ra sao.
Nếu cần một trạm đo độc lập, bộ ghi dữ liệu kiểu CR1000X sẽ tiện vì nó tích hợp
sẵn mọi thứ. Nếu nhà máy đã dùng PLC Modicon X80 rồi, mô-đun analog tương thích
sẽ dễ bảo trì hơn. Còn nếu tự thiết kế bo mạch nhúng, IC RS-485 chỉ là phần giao
tiếp, không thay thế được bộ thu thập dữ liệu.

\section{Kết quả và kỹ năng thu nhận}

Sau quá trình tìm hiểu, em đã hệ thống hóa được bốn nhóm cung ứng của TECHPAL,
vẽ được sơ đồ liên hệ giữa Sale, Mua hàng Quốc tế và Logistics, lập một bảng
kiểm kỹ thuật và đối chiếu bốn mã sản phẩm với datasheet của hãng. Bài toán nhiệt
độ giúp em thử áp dụng cách chuyển từ yêu cầu chức năng sang kiến trúc và danh
sách thông tin cần kiểm tra.

Điều em thấy quan trọng nhất là không nên ra quyết định chỉ vì tên sản phẩm nghe
hợp lý hay mô tả ngắn trên web trùng khớp. Mỗi mã linh kiện đều cần kiểm tra
bằng datasheet, tài liệu hệ thống và trạng thái vòng đời. Ngoài kiến thức chuyên
ngành, công việc cung ứng còn đòi hỏi khả năng đọc tiếng Anh, ghi chép cẩn thận
và phối hợp giữa nhiều bộ phận. Báo cáo của em dừng ở mức khảo sát tài liệu và
đề xuất phương án, chưa có nguyên mẫu hoặc số liệu đo thực tế.
